\section{Interface bridge assumptions and failure modes (one-page audit map)}
\label{app:interface_bridge_assumptions_failure_modes}

\paragraph{Purpose.}
\AuditTag This appendix is a reader-facing audit map.
It collects the most important \InterfaceTag/\MatchTag bridge assumptions that connect finite protocol objects to physical-language statements (energy scales, horizons, and CP-odd normalizations), and it records concrete failure modes and falsification routes.
No theorem-level folding statement depends on this appendix as a premise.

\paragraph{How to read this table.}
Each row has four parts: (i) the bridge assumption, (ii) where it is used, (iii) what would falsify it (or force revision), and (iv) what remains intact if it fails (the minimal core that stays true).

\begin{table}[H]
\centering
\scriptsize
\setlength{\tabcolsep}{6pt}
\renewcommand{\arraystretch}{1.2}
\begin{tabular}{p{0.22\textwidth} p{0.23\textwidth} p{0.25\textwidth} p{0.22\textwidth}}
\toprule
bridge assumption (type) & used in & falsification / failure mode & minimal core that remains \\
\midrule
\textbf{(B1) Staircase calibration} \MatchTag\ $\mu_{\mathrm{th}}(m)=m_e\varphi^{r_{\mathrm{step}}(m-6)}$ with $r_{\mathrm{step}}=2\pi$ (bounded family) &
Table~\ref{tab:resolution_thresholds}, Table~\ref{tab:resolution_spectrum}, Table~\ref{tab:m_energy_dictionary_compact}; calibration-dependent prediction P3 (Section~\ref{subsec:calibration_dependent_predictions}) &
If threshold onsets in data systematically avoid the calibrated bands without an alternative $r_{\mathrm{step}}$ inside the stated bounded family; or if the chosen anchors $(m=6\leftrightarrow m_e)$ fail under updated conventions &
All calibration-free protocol predictions (P1/P2/P4/P5 etc.) and all finite folding statements; only the $\mu$-locations of thresholds must be re-calibrated \\
\midrule
\textbf{(B2) Protocol capacity} \InterfaceTag\ $I_{\mathrm{prot}}(m,n;R)=m|R|$ (bits) &
Holographic capacity bound (Appendix~\ref{app:holo_capacity_bound}); budget-triggered horizon rules and occupancy audits (Appendix~\ref{app:protocol_horizon_tick_trap}; Section~\ref{sec:qg_interface_full_fusion}) &
If the per-site register size is not $2^m$ in the declared protocol class (e.g.\ analog degrees of freedom not captured by $m$ bits, or a non-tensor boundary algebra) &
Finite folding core and combinatorics at fixed $m$; capacity-facing horizon claims must be reformulated under the revised boundary model \\
\midrule
\textbf{(B3) High-$m$ $\Rightarrow$ horizon-like trapping (interface)} &
Section~\ref{subsec:bh_high_m_particles}; full-fusion horizon/trap/leakage modules (Section~\ref{sec:qg_interface_full_fusion}) &
If increasing $m$ (or observer budget) does not induce saturation/occupancy regimes in protocol audits, or if delay/leakage proxies do not exhibit the predicted scaling/transition at the same boundary &
All capacity-only occupancy and entropy bounds remain valid; only the physical interpretation as ``BH-like'' must be downgraded \\
\midrule
\textbf{(B4) CP-odd multiplicity} \InterfaceTag\ $d_{\mathrm{CP}}=\dim\mathfrak{su}(3)+\dim\mathfrak{su}(2)$ under rephasing invariance and explicit projective quotients &
The $J$ normalization dictionary (Section~\ref{subsec:cp_jarlskog}; Proposition~\ref{prop:dcp_forced_11}) &
If CP-odd residue requires an additional independent abelian contribution not already represented in the phase-space quotient, or if gauge closure differs from $SU(3)\times SU(2)\times U(1)$ within the declared interface contract &
The $\pi^7$ phase-space factorization and the audit tables remain; only the discrete multiplicity dictionary (the ``11'') must be modified consistently \\
\midrule
\textbf{(B5) HTF-lite trace bridge (audit-only)} \AuditTag\ Assumption~\ref{ass:htf_lite_operator_pack} &
Interpretive unification (Section~\ref{subsec:interpretive_unification_complex_exp}); operator mother space notes (Appendix~\ref{app:operator_mother_space}) &
If no declared kernel family yields a meromorphic identity with resolvent-mode factors whose principal parts are preserved under the stated renormalization; or if detectability/noncancellation fails (interior poles are removed by same-point cancellations) &
All theorem-level folding statements; all protocol-level wormhole modules and ledgers (Definition~\ref{def:wormhole_pointer_jump} and Section~\ref{sec:qg_interface_full_fusion}); the pole-barrier template remains usable as a generic analyticity gate (Appendix~\ref{app:trace_pole_barrier_template}) \\
\midrule
\textbf{(B6) Local interface factorization gate (Woodbury/Schur)} \AuditTag\ (LF1) &
Equivalence-ladder registry and operator-mother-space update bookkeeping (Appendix~\ref{app:equivalence_ladder_certificates}; Appendix~\ref{app:operator_mother_space}); any shortcut/update claim promoted to det/logdet/delay &
If an update/shortcut claim does not provide an explicit baseline/counterfactual split and a small-dimensional update factor (finite-rank Woodbury or dimension-changing Schur complement), or if the factor depends on undeclared degrees of freedom &
Downgrade to paid ledger + counterfactual baseline only (no determinant/logdet promotion); keep theorem-level folding chain unchanged; treat remaining statements as audit narrative only \\
\bottomrule
\end{tabular}
\caption{Bridge assumptions, failure modes, and minimal-core fallbacks. This table is an audit-facing interface map, not a theorem source.}
\label{tab:bridge_assumptions_failure_modes}
\end{table}

